\subsection{Same-Oracle Optimistic Support and Neural-Action Attribution}
\label{sec:r6_bank_evidence}
Table \ref{tab:r6_ols} compares two coverage objectives with an identical exact Bellman oracle, action menu, state space, and tolerance. The initial-distribution version concentrates all weight on $(u,X)=(2,1.25)$ at date zero. The statewise version takes the maximum gap across all 1,617 states and all eight decision dates. Both start from $d=0,1$, insert a worst remaining corner, and use independently evaluated policy features. We report common kernel construction separately; bank memory counts policies and features, excluding shared transition arrays.
\begin{table}[!htbp]
\centering\small
\caption{Exact-Oracle Scalar Optimistic Support}
\label{tab:r6_ols}
\begin{tabular}{lrrrr}
\toprule
Coverage objective & Anchors & Seconds & MiB & Global bound \\
\midrule
Initial state/date & 3 & 5.82463 & 1.14732 & 0.0529688 \\
All states/dates & 47 & 91.2708 & 17.9746 & 0.000930698 \\
\bottomrule
\end{tabular}
\par\smallskip
\begin{minipage}{0.98\linewidth}\footnotesize
Notes: Same exact full-menu oracle and $10^{-3}$ target. The initial-state bound is 0.000410746. Global means all finite states, dates, and $d\in[0,1]$. Common kernel setup costs 2.78081 seconds and is excluded from both preparation columns. MiB counts policies and features, not shared kernels or upper buffers. The independent full-line/adjacent-line discrepancy is 4.44089e-16.
\end{minipage}
\end{table}

The initial-distribution target needs only three anchors. Its much larger all-state/date bound demonstrates why that guarantee cannot be silently used for a different population or a later date. The statewise calculation uses 47 anchors. An independent calculation with every policy line reproduces its bound and verifies the adjacent-policy identity to floating-point tolerance. This is an executed comparison of guarantee scope and the scalar-OLS correspondence, not evidence that a new name reduces the oracle count.

We then remove every neural proposal from the deployed lower policies, leaving the same 1,565-action common mesh. The upper information continues to come from the original full, neural-inclusive target. Table \ref{tab:r6_mesh} reports the anchor losses and the maximum over every contract interval. The mesh-only bank remains within $10^{-3}$ of the \emph{original} target for all $d\in[0,1]$, states, and dates. This is not a comparison against an easier restricted optimum, and it is not an autonomous mesh-only anchor-discovery experiment: the 47 locations and the full-target upper oracle are shared.
\begin{table}[!htbp]
\centering\small
\caption{Removing Neural Actions from the Feasible Policy Bank}
\label{tab:r6_mesh}
\begin{tabular}{lrrr}
\toprule
Feasible lower policies & Anchors & Anchor loss & Interval bound \\
\midrule
Full menu & 47 & 0 & 0.000930698 \\
Common mesh only & 47 & 3.87125e-05 & 0.000930698 \\
\bottomrule
\end{tabular}
\par\smallskip
\begin{minipage}{0.98\linewidth}\footnotesize
Notes: Both rows use the original neural-inclusive optimal upper target at the same anchors. The second row removes all three neural candidates only from the lower policies. The interval bound covers all states, dates, and $d\in[0,1]$. Mesh-policy construction costs 88.0996 seconds; this is not an autonomous mesh-anchor-discovery time.
\end{minipage}
\end{table}

The neural actions have a small, nonzero value contribution at some nodes, but are unnecessary for this tolerance. Exact dynamic programming constructs the repaired bank; independent policy evaluation supplies its features; the upper oracle supplies the welfare benchmark; policy switching supplies deployment. The initial trained proposals and their profitable deviations remain reported. The transition-law theorem is an additional guarantee applicable to either neural or mesh policies; it does not require reassigning these observed roles to learning.

\subsection{A Reusable Structural Solver for the Resource Economy}
\label{sec:structural_qp}
Stack the $21$ nonanticipative node controls into $z\in\R^{21d}$. For every initial state $x$, the same model has cost
\begin{equation}
 C_x(z)=\tfrac12z^{\mathsf T}Hz+(Fx+f_0)^{\mathsf T}z+c(x),
 \qquad z_h\geq0,\quad\mathbf1^{\mathsf T}z_h\leq0.25d.
 \label{eq:reusable_qp}
\end{equation}
The matrices $H,F$ and vector $f_0$ are assembled once from the affine state paths. We scale the controls by their node probabilities, compute the resulting Hessian's largest eigenvalue once, and apply accelerated projected minimization. The scaling is constant within a node, so the node's feasible-set projection is unchanged. Each query starts at zero controls. The computation uses no future-shock information, stored query answer, or cross-query warm start.

The comparison follows the structural counterexample in the fifth-round advisory report, with an explicit independent assembly in the revision source. We check the assembled costs and gradients against the original full-tree routine. Every query is evaluated on all 64 terminal paths, and its stopping gap is recomputed from the original tree gradient, not just the assembled matrix. Table \ref{tab:r6_qp} uses the same fresh initial states and $10^{-3}$ target as the earlier workload. Setup is included for the QP; critic-only retraining, zero-start improvement, and full policy evaluation are included for the learned arm. Query timings are medians of three serial single-thread runs.
\begin{table}[!htbp]
\centering\small
\caption{Fresh Queries with Reusable Quadratic Structure}
\label{tab:r6_qp}
\begin{tabular}{lrrrr}
\toprule
Queries & Learned total & QP total & QP query & QP max gap \\
\midrule
32 & 0.884792 & 0.00759194 & 0.00434819 & 0.000633261 \\
128 & 0.964309 & 0.014003 & 0.0107592 & 0.000664162 \\
512 & 1.57165 & 0.0405753 & 0.0373315 & 0.000664162 \\
2048 & 3.66748 & 0.143248 & 0.140004 & 0.000677426 \\
\bottomrule
\end{tabular}
\par\smallskip
\begin{minipage}{0.98\linewidth}\footnotesize
Notes: Times are seconds, single-thread, query medians of three. The learned total includes critic retraining and full policy evaluation; QP total includes structural setup and independent full-tree gradient/gap evaluation. All queries start at zero controls. The largest learned policy-loss upper bound is 0.000675399. The original reference timings and all timing samples remain in the JSON record.
\end{minipage}
\end{table}

The reusable QP is faster in these workloads, including online time alone. The complete-tree gaps certify its costs at the same tolerance. This is a substantive change in the computational conclusion, not a revision of the historical elapsed times. A fitted PSD-quadratic \emph{continuation approximation} and the exact parametric quadratic \emph{program} are different comparators. The recorded failure of the former does not imply a failure of the latter. Consequently this resource model supports controlled neural approximation and component attribution, but supplies no empirical basis for choosing learning over the reusable structural solver. The additional transferable guarantee studied next concerns a nonlinear stopped economy with changed transition laws, not a retiming of this quadratic example.

\subsection{A Continuum of Shock-Law Counterfactuals}
\label{sec:r6_transition_experiment}
Let law 0 be the original one-step quadrature with shock correlation $-0.25$, and law 1 the otherwise identical quadrature with correlation $0.25$. At every date, law 1 occurs with probability $\lambda$, independently across dates and without being known before that date's control. The state transition, exit payment, and discounted continuation are the corresponding mixtures. The state grid, consumption and portfolio bounds, adjustment menu, terminal payoff, and three frozen neural proposals remain unchanged. This is a change in the frequency of positive versus negative co-movement regimes in the finite economy. It is not the claim that interpolating two quadratures equals discretizing every intermediate Brownian correlation.

For each of the contract contrasts $d=0,0.5,1$, we solve endpoint problems and adaptively subdivide the transition-parameter interval using Theorem \ref{thm:corrected_chord}. The endpoint policies are evaluated by \eqref{eq:bernstein_recursion}. Eight closed Bernstein cells within each anchor interval bound its full policy-value gap; no target problem is solved at those cell endpoints merely to label a grid uniform. Table \ref{tab:r6_kernel} reports the maximum certificate over every state, date, and $\lambda\in[0,1]$. The exact parameter remains common across the horizon. The certificate is distinct from the older reward interval: the displayed $d$ values are three fixed contracts, not a new joint continuum claim over $(d,\lambda)$.
\begin{table}[!htbp]
\centering\small
\caption{Uniform Policy Reuse across Transition Mixtures}
\label{tab:r6_kernel}
\begin{tabular}{lrrrr}
\toprule
$d$ & Anchors & Corrected bound & Count bound & Prep. seconds \\
\midrule
0 & 19 & 0.000904573 & 0.0127756 & 52.5879 \\
0.5 & 18 & 0.000899829 & 0.0123306 & 58.5677 \\
1 & 17 & 0.000941412 & 0.00928855 & 53.1281 \\
\bottomrule
\end{tabular}
\par\smallskip
\begin{minipage}{0.98\linewidth}\footnotesize
Notes: Each row covers every $\lambda\in[0,1]$, all 1,617 states, and eight decision dates at the displayed fixed contract. Corrected preparation includes endpoint solves, policy-coefficient construction, and interval certification, excluding the common two-kernel build (5.43095 seconds) and subsequent independent replay. The count bound uses its separate informed upper oracle and three feasible anchor policies. It is valid but does not meet the $10^{-3}$ target. No joint $(d,\lambda)$-continuum or diffusion guarantee is asserted.
\end{minipage}
\end{table}

For comparison, a separate count-informed upper recursion, coupled with the three feasible policies at $\lambda=0,0.5,1$, gives a valid but looser uniform bound. Its gap includes the value of advance count information. Reporting it beside the corrected-chord result makes clear that an upper-bound method can be valid without being sufficiently sharp for the desired welfare decision. The corrected construction adds solves where its actual bound requires them and stops only when the $10^{-3}$ target is attained.

Every final anchor interval is independently replayed at its endpoints and midpoint using a complete backward optimization of the same mixture model. Fixed-policy polynomial values are also checked against a direct selected-kernel recursion. Small independent models enumerate all regime sequences, and two-date models enumerate every deterministic Markov policy. These tests distinguish a correct polynomial identity from an accidentally favorable main-model result. The uncorrected-chord negative control has endpoint values zero and midpoint value $1/4$; the upper correction is therefore a necessary substantive component, not additional bookkeeping around an already valid reward chord.

All calculations use double precision. The transition arrays are positive, the prescribed finite actions are exhaustively checked, and recorded arithmetic discrepancies are compared with explicit tolerances. These are finite-model numerical certificates, not directed-rounding interval arithmetic or a diffusion discretization guarantee. A CSR representation accelerates multiplication by the identical transition rows; its values are checked against the inherited array implementation, while fixed-policy evaluation retains a separate arithmetic path.

\subsection{The Preference-Formation Contribution to Risk Choice}
\label{sec:r6_mechanism_evidence}
Table \ref{tab:r6_option} evaluates the risk-class comparison in Section \ref{sec:risk_frontier}. Both arms retain preference shocks, the original covariance, the same operating boundary and settlement, and the full consumption and portfolio menus. The restricted arm sets deliberate adjustment to zero at \emph{all} states and dates. The difference $\Omega_+-\Omega_-$ is computed from separately optimized and independently evaluated risk-class policies, not from a first-order condition imposed on a binding control.
\begin{table}[!htbp]
\centering\small
\caption{Deliberate Preference Formation and Risk-Class Values}
\label{tab:r6_option}
\begin{tabular}{lrrr}
\toprule
$d$ & $\Delta^{\rm adj}$ & $\Delta^0$ & $\Omega_+-\Omega_-$ \\
\midrule
0 & -0.000875905 & -0.00133454 & 0.000458639 \\
0.25 & -0.000276562 & -0.000730716 & 0.000454155 \\
0.35 & -4.62573e-05 & -0.000489849 & 0.000443592 \\
0.365 & -1.17117e-05 & -0.000453738 & 0.000442026 \\
0.375 & 1.13188e-05 & -0.000429663 & 0.000440982 \\
0.4 & 6.88949e-05 & -0.000369477 & 0.000438372 \\
0.5 & 0.000299199 & -0.000135684 & 0.000434883 \\
1 & 0.00128126 & 0.000950467 & 0.000330789 \\
\bottomrule
\end{tabular}
\par\smallskip
\begin{minipage}{0.98\linewidth}\footnotesize
Notes: Positive $\Delta$ favors a positive first risky position; negative $\Delta$ favors the nonpositive class. Both classes optimize every subsequent action. The restricted arm sets deliberate adjustment to zero at all states and dates but retains preference shocks and every financial and stopping primitive. Each reported value is independently evaluated in the same cardinal utility units.
\end{minipage}
\end{table}

At $d=0.5$, the adjustment-enabled economy favors a positive first risky position, whereas the otherwise identical economy without deliberate adjustment favors the nonpositive class. The same disagreement occurs at $d=0.375$ and $d=0.4$. Thus the relative preference-formation option is large enough to reverse the risk ranking under an unchanged stopping contract. At $d=1$, both economies favor the positive class: the operating-contract incentive can produce a reversal without deliberate preference formation, but the formation option changes where that reversal occurs. This separates a generic duration incentive from a preference-specific ranking effect through \eqref{eq:option_decomposition}. Neither fixed consumption at a bound nor movement between endpoint portfolios alone supplies this identification.

Table \ref{tab:r6_frontier} holds the financial and stopping primitives fixed and varies the adjustment-cost coefficient. It reports opposite-sign brackets and the attaining policies' duration and effort differences. The ratio in \eqref{eq:risk_frontier_slope} is a local regular-branch prediction in cardinal flow-utility units per unit of adjustment cost, not a global derivative through every policy kink. We do not claim uniqueness from a bisection bracket; uniqueness requires the uniform duration-separation condition of Proposition \ref{prop:risk_frontier}. The no-adjustment bracket separately shows which part of the risk response remains when the formation option is removed.
\begin{table}[!htbp]
\centering\small
\caption{Risk Indifference and the Duration--Effort Comparison}
\label{tab:r6_frontier}
\begin{tabular}{lrrrr}
\toprule
$k$ / restriction & Opposite-sign bracket & $A_+-A_-$ & $C_+-C_-$ & Ratio \\
\midrule
0.5 & [0.34687805, 0.34689331] & 0.00230304 & 4.6285e-05 & 0.0200973 \\
2 & [0.37007141, 0.37008667] & 0.00230304 & 2.12312e-05 & 0.00921875 \\
8 & [0.48182678, 0.48184204] & 0.00230304 & 7.90303e-05 & 0.0343156 \\
No adjustment & [0.55891418, 0.55892944] & 0.00230304 & 0 & 0 \\
\bottomrule
\end{tabular}
\par\smallskip
\begin{minipage}{0.98\linewidth}\footnotesize
Notes: Each bracket has opposite risk-advantage signs at its endpoints and width at most $2\times10^{-5}$. It proves existence of a crossing, not global uniqueness. Features and their ratio are evaluated at the midpoint. The ratio predicts a regular local branch slope only under the stated nonzero-duration and differentiability conditions. A zero-adjustment policy has zero adjustment effort. Bracket endpoints and their signed values are stored, rather than inferred from rounded table entries.
\end{minipage}
\end{table}

Increasing the adjustment-cost coefficient moves the observed opposite-sign bracket toward higher duration compensation. The no-adjustment bracket lies further to the right. These results accord with the positive effort--duration ratios of the attaining class policies, subject to the local regularity conditions. The duration--effort and option-value propositions state the portable mechanism; the matched finite economy identifies its effect on an actual risk ranking. Changing the operating boundary or settlement curvature would define further economic contracts, rather than innocuous numerical robustness checks. Their effects are not identified by the present matched restriction. The original boundary calculation, correlation panel, and joint numerical refinements remain available and are not silently replaced by the new experiment.
